\paragraph{Retail, e-commerce, and supply-chain decision benchmarks.}
Prior work has studied retail, e-commerce, and supply-chain decision-making through benchmarks for retail simulation and promotion optimization~\citep{xia2023retailsynth,xia2024simulationretailpromotions}, demand recovery and transaction modeling~\citep{wang2025freshretailnet,tkachuk2024consumertransactions}, web-shopping agents~\citep{wang2025shoppingbench,3webshopscalablerealworldweb}, and supply-chain or inventory control~\citep{OFCOURSE,inventory,leluc2023marlimmultiagentreinforcementlearning,invagentlargelanguagemodel,aimbenchevaluatingdecisionmakingbiases}. These benchmarks cover important components of retail and supply-chain operation, but typically isolate subproblems such as promotion, demand modeling, shopping interaction, or inventory control. RetailBench instead evaluates tool-using LLM agents in a coupled retail environment where pricing, replenishment, supplier choice, shelf assortment, customer feedback, exogenous events, and financial constraints interact over time.

\paragraph{Long-horizon agent benchmarks.}
Recent benchmarks increasingly evaluate LLM agents beyond isolated task completion, including ultra-long-horizon interaction, realistic office workflows, virtual-world planning, lifelong experience reuse, workplace task execution, and long-running business operation~\citep{utltabench,odysseybenchevaluatingllmagents,herobenchbenchmarklonghorizonplanning,zheng2025lifelongagentbench,xu2024theagentcompany,backlund2025vendingbenchbenchmarklongtermcoherence,andonlabs2025vendingbench2}. These benchmarks highlight the importance of persistence, memory, exploration, and workflow-level decision making. However, they generally do not model data-grounded retail dynamics in which pricing, replenishment, supplier choice, shelf visibility, customer feedback, external events, and cash flow co-evolve over many days. RetailBench complements this line of work by grounding long-horizon agent evaluation in closed-loop product-level retail operation and by supporting diagnostic analysis of evidence acquisition, evidence-to-action conversion, and long-horizon policy consistency.